\documentclass[a4paper,12pt,oneside]{book}
\linespread{1.5}
\usepackage{geometry}
\geometry{top=2cm,bottom=2.5cm,left=3.5cm,right=3.5cm} % Impostiamo il valore dei margini.
\usepackage[utf8]{inputenc}
\usepackage{tikz,pgf}
\usepackage{amsfonts}
\usepackage[italian]{babel} %% Abbiamo istruito LaTeX ad usare la lingua italiana.
\usepackage{todonotes}
\usepackage{subfig}
\usepackage{fancyhdr}
\setlength{\parindent}{0cm}

\usepackage{graphicx}
\usepackage{float}
\usepackage{csquotes}
\usepackage{amsmath}
\usepackage{numprint}
\usepackage{pgfplots}
\usepackage{pgfplotstable}
\usepackage{tocbibind}
\setcounter{tocdepth}{3}
\usepackage{hyperref}

\usepackage{fancyhdr}
\pagestyle{fancy}
\renewcommand{\headrulewidth}{0pt}
\fancyhf{}
\fancyhead[L]{}
\fancyfoot[C]{\thepage}

\graphicspath{{imgs/}} %% cartella dalla quale ricavare le immagini %%

\setlength {\marginparwidth }{2cm} 
\pgfplotsset{compat=1.16}

\newenvironment{dedication}
  {\clearpage           % we want a new page
   \thispagestyle{empty}% no header and footer
   \vspace*{\stretch{1}}% some space at the top 
   \itshape             % the text is in italics
   \raggedleft          % flush to the right margin
  }
   {\par % end the paragraph
   \vspace{\stretch{3}} % space at bottom is three times that at the top
   \clearpage           % finish off the page
  }
  
\begin{document}
%% Frontespizio %%
\begin{titlepage}
    \begin{center}
    \vspace*{0.2cm}
    \begin{figure}[H]
        \centering
        \includegraphics[width=0.8\textwidth]{imgs/logo.png}
    \end{figure}
        \vspace{5mm}
%% Dipartimento %%
    	\uppercase{\normalsize Dipartimento di Ingegneria Informatica, Modellistica, Elettronica e Sistemistica}\\
    \end{center}
    \begin{center}
%% Corso di Laurea %%
    	\large{ Corso di Laurea (inserire corso)}\\
        \vspace{3mm}
        \large{\bf Tesi di Laurea}\\
    	\vspace{10mm}
%% Titolo della Tesi %%
    {\LARGE{\bf Titolo Tesi (inserire titolo)}}\\	
     \vspace{7mm}
    \end{center}
    
    \vspace{25mm}
    \noindent
    \begin{minipage}[t]{0.47\textwidth}
%% Relatore %%
    	{\large{ Relatore:\vspace*{0.22cm}\\\bf Prof. \\Nome Cognome}}
    	\vspace{12mm}
    \end{minipage}
    \hfill
    \begin{minipage}[t]{0.4\textwidth}\raggedleft
%% Candidato %%
    	{\large{Candidato: \\ \bf Nome Cognome\\ Mat. 123456}}
    \end{minipage}
    
    \vspace{30mm}
    
%% Anno Accademico %%
    \centering{\large \uppercase{ Anno Accademico 12/34}}

\end{titlepage}

\begin{dedication}
Inserire dedica e ringraziamenti...
\end{dedication}

%% Corpo della tesi %%

%%\chapter*{Sommario} %% In questa zona inseriamo l'abstract del nostro documento. %%
%%Inserisci abstract: ...


\tableofcontents
\input{capitoli/capitolo1}
\input{capitoli/capitolo2}
\input{capitoli/capitolo3}
\input{capitoli/capitolo4}
\input{capitoli/capitolo5}
\input{capitoli/capitolo6}
\input{capitoli/capitolo7}

\begin{thebibliography}{100} %% 100 rappresenta il n° massimo di voci all'interno della bibliografia: può essere modificato a vostro piacimento. %%
\bibitem{[1]:}
Inserire qui la bibliografia...
\bibitem{[2]:}
Inserire qui la bibliografia...

\end{thebibliography}
\end{document}
